\begin{trapbox}
    \begin{description}[style=nextline, leftmargin=2.8em, labelsep=0.5em, itemsep=0.36em]
      \item[\textbf{\textcolor{CTrap}{易错点一（换元后忘记新变量的范围）}}]
        把 $u=t+1/t$ 当作任意实数，解出 $u=1$ 后直接回代。
      \item[\textbf{\textcolor{CTrap}{正确理解（先写 $|u|\ge2$）}}]
        对非零实数 $t$，$u^2-4=(t-1/t)^2\ge0$；取 $u=1$ 会产生方程 $t^2-t+1=0$，无实根。
      \item[\textbf{\textcolor{CTrap}{错因分析}}]
        整体换元是把原变量映射到一个更小的取值集合；如果扩大新变量的范围，会出现不属于原题的假解。
    \end{description}

    \medskip
    \begin{description}[style=nextline, leftmargin=2.8em, labelsep=0.5em, itemsep=0.36em]
      \item[\textbf{\textcolor{CTrap}{易错点二（把倒数和与倒数差的取值范围混用）}}]
        令 $v=t-1/t$ 后，也写成 $|v|\ge2$。
      \item[\textbf{\textcolor{CTrap}{正确理解（差型可以取零）}}]
        \begin{center}
\begin{tikzpicture}[font=\small]
\begin{axis}[at={(0,0)},width=6.3cm,height=6cm,axis lines=middle,
 xmin=-3.5,xmax=3.5,ymin=-4.5,ymax=4.5,xtick={-1,1},ytick={-2,2},
 yticklabel style={anchor=west,xshift=4pt},
 xlabel={$t$},title={和型 $u=t+t^{-1}$}]
\addplot[CStatement,very thick,domain=.24:3.4,samples=130] {x+1/x};
\addplot[CStatement,very thick,domain=-3.4:-.24,samples=130] {x+1/x};
\addplot[only marks,CStatement] coordinates {(1,2) (-1,-2)};
\end{axis}
\begin{axis}[at={(6.5cm,0)},width=6.3cm,height=6cm,axis lines=middle,
 xmin=-3.5,xmax=3.5,ymin=-4.5,ymax=4.5,xtick=\empty,ytick=\empty,
 xlabel={$t$},title={差型 $v=t-t^{-1}$}]
\addplot[CExample,very thick,domain=.22:3.4,samples=130] {x-1/x};
\addplot[CExample,very thick,domain=-3.4:-.22,samples=130] {x-1/x};
\addplot[only marks,CExample] coordinates {(1,0) (-1,0)};
\end{axis}
\end{tikzpicture}
\end{center}

        看同尺度曲线：和型缺少中间纵坐标，差型却经过 $(1,0)$、$(-1,0)$。任意实数 $v$ 都可取得，理由见证明补充。
      \item[\textbf{\textcolor{CTrap}{错因分析}}]
        $t+1/t$ 的两项同号，$t-1/t$ 的两项可能抵消；不能套用相同的基本不等式条件。
    \end{description}

    \medskip
    \begin{description}[style=nextline, leftmargin=2.8em, labelsep=0.5em, itemsep=0.36em]
      \item[\textbf{\textcolor{CTrap}{易错点三（错误的高次展开）}}]
        写成 $t^2+t^{-2}=u^2$，或写成 $t^3+t^{-3}=u^3$。
      \item[\textbf{\textcolor{CTrap}{正确理解（交叉项不能漏）}}]
        \[
          t^2+t^{-2}=u^2-2,\quad t^3+t^{-3}=u^3-3u.
        \]
        例如 $t=1$ 时，左边均为 $2$，而 $u=2$；$u^2=4$、$u^3=8$ 显然不等于 $2$。
      \item[\textbf{\textcolor{CTrap}{错因分析}}]
        将 $(a+b)^n$ 误当成 $a^n+b^n$。用 $t=1$ 检验是最小反例。
    \end{description}

    \medskip
    \begin{description}[style=nextline, leftmargin=2.8em, labelsep=0.5em, itemsep=0.36em]
      \item[\textbf{\textcolor{CTrap}{易错点四（看到首尾系数一样，就断定一定可换）}}]
        只看最高次和常数项相等，没有检查中间各项成对的系数。
      \item[\textbf{\textcolor{CTrap}{正确理解（完整系数都要成对相同）}}]
        \begin{center}
\begin{tikzpicture}[x=1cm,y=.75cm,font=\small]
\node[font=\bfseries] at (5.5,5.6) {首尾相等还不够：中间各对也要相等};
\foreach \x/\k in {1/4,3.25/3,5.5/2,7.75/1,10/0} {\node at (\x,4.65) {$t^{\k}$};}
\foreach \x/\v in {1/1,3.25/2,5.5/3,7.75/2,10/1} {
\filldraw[fill=CExample!18,draw=white] (\x-.4,2.7) rectangle +(0.8,\v*.4);
\node at (\x,2.4) {$\v$};}
\foreach \x/\v in {1/1,3.25/2,5.5/3,7.75/5,10/1} {
\filldraw[fill=CTrap!20,draw=white] (\x-.4,-.3) rectangle +(0.8,\v*.4);
\node at (\x,-.65) {$\v$};}
\draw[CExample] (.4,2.7)--(10.6,2.7);\draw[CTrap] (.4,-.3)--(10.6,-.3);
\node[anchor=east,text=CExample] at (.1,3.4) {可配对};
\node[anchor=east,text=CTrap] at (.1,.5) {不对称};
\node at (5.5,1.8) {上行：$2t+2/t=2u$};
\node at (5.5,-1.35) {下行：$2t+5/t$ 的两个系数不等，不能直接合成 $au$。};
\end{tikzpicture}
\end{center}

        同一高度单位表示同一系数大小；下行虽首尾都是 $1$，但 $t^3$ 与 $t$ 的系数 $2,5$ 不同，配对失败。
      \item[\textbf{\textcolor{CTrap}{错因分析}}]
        除以 $t^2$ 后，有 $2t+5/t$，不是 $2(t+1/t)$；只匹配首尾不足以保证整体配对。
    \end{description}

    \medskip
    \begin{description}[style=nextline, leftmargin=2.8em, labelsep=0.5em, itemsep=0.36em]
      \item[\textbf{\textcolor{CTrap}{易错点五（仅凭倒数和决定倒数差的符号）}}]
        看到 $t^2-t^{-2}$，用 $u=t+1/t$ 后直接把它写成 $u\sqrt{u^2-4}$。
      \item[\textbf{\textcolor{CTrap}{正确理解（平方可以，符号要单独判断）}}]
        \[
          (t^2-t^{-2})^2=u^2(u^2-4),
        \]
        但 $t^2-t^{-2}=u(t-1/t)$，其中 $t-1/t$ 的符号随 $t$ 改变。理解页两根虚线段给出了 $t=2$ 与 $t=1/2$ 的反号例子。
      \item[\textbf{\textcolor{CTrap}{错因分析}}]
        同一个 $u$ 可能对应 $t$ 与 $1/t$ 两个不同的原变量；平方后符号信息被消去，开平方不能擅自取正。
    \end{description}

    \medskip
    \begin{description}[style=nextline, leftmargin=2.8em, labelsep=0.5em, itemsep=0.36em]
      \item[\textbf{\textcolor{CTrap}{易错点六（除以 $t^n$ 前没检查 $t=0$）}}]
        对任何高次方程直接除以 $t^n$，可能漏掉 $t=0$ 的解。
      \item[\textbf{\textcolor{CTrap}{正确理解（先验证零根）}}]
        例如 $t^5+t^3=0$ 有根 $t=0$，不能未经检查就除以 $t^3$。本节四次首尾对称例子常数项为 $1$，所以 $t=0$ 本来就不可能是根。
      \item[\textbf{\textcolor{CTrap}{错因分析}}]
        除以变量的幂等于排除了该变量为零的情况，必须先验证是否会丢解。
    \end{description}

    \medskip
    \begin{description}[style=nextline, leftmargin=2.8em, labelsep=0.5em, itemsep=0.36em]
      \item[\textbf{\textcolor{CTrap}{易错点七（只求出换元后的最值，没有检验能否取得）}}]
        把 $z=t^2+t^{-2}$ 当作任意实数，对 $z^2-3z-2$ 使用顶点 $z=3/2$ 求最小值。
      \item[\textbf{\textcolor{CTrap}{正确理解（顶点不在允许区间）}}]
        $z\ge2$，故最小值在 $z=2$，为 $-4$，由 $t=\pm1$ 取得。不能用不在取值范围内的顶点。
      \item[\textbf{\textcolor{CTrap}{错因分析}}]
        降次只是改变式子的外观，不会扩大新变量实际允许的范围。
    \end{description}
\end{trapbox}
